\section{Find How Many Times an Array Has Been Rotated}
\label{sec:bs-how-many-rotations}

\problemstatement{Given a sorted array that has been rotated some number of
times (all elements distinct), determine \textbf{how many times} it has
been rotated -- equivalently, find the number of left rotations applied to
the original sorted array.}

\begin{patternbox}
This problem is solved by noticing that it is asking for almost the same
quantity as Section~\ref{sec:bs-find-min-rotated}, phrased differently: the
number of left rotations applied to a sorted array equals exactly the
\textbf{index} of the minimum element. If the original sorted array is
rotated left by $k$ positions, the element that used to be at index $0$
(the global minimum) ends up sitting at index $k$. Whenever a problem
about a rotated array asks for ``how many rotations,'' translate it
immediately into ``find the index of the minimum.''
\end{patternbox}

\subsection*{Understanding the problem carefully}

Picture the original sorted array $[v_0 < v_1 < \cdots < v_{n-1}]$. Rotating
it left by $k$ positions produces
$[v_k, v_{k+1}, \dots, v_{n-1}, v_0, v_1, \dots, v_{k-1}]$. The smallest
value $v_0$ now sits at position $k$ in the rotated array. So finding the
number of rotations is exactly finding the position of the minimum element
-- the same quantity computed in Section~\ref{sec:bs-find-min-rotated},
except we now want the \emph{index} rather than the \emph{value}.

\begin{ideabox}[Reuse, Don't Re-derive]
Apply the exact same binary search as
Section~\ref{sec:bs-find-min-rotated} to locate the minimum, but track and
return $\textit{lo}$ (the index) instead of $\textit{arr}[\textit{lo}]$
(the value).
\end{ideabox}

\subsection*{Algorithm}

Identical to Algorithm~\ref{sec:bs-find-min-rotated}'s \textsc{FindMin}
procedure, returning $\textit{lo}$ instead of $\textit{arr}[\textit{lo}]$.

\subsection*{Dry run}

\dryrun{Take $\textit{arr} = [4,5,6,7,0,1,2]$ (same array as
Section~\ref{sec:bs-find-min-rotated}).}

The binary search trace is identical to before, terminating with
$\textit{lo} = 4$. So the array has been rotated $4$ times -- indeed, the
original sorted array would have been $[0,1,2,4,5,6,7]$, and rotating it
left by $4$ positions gives exactly $[4,5,6,7,0,1,2]$.

\subsection*{C++ implementation}

\begin{lstlisting}
#include <bits/stdc++.h>
using namespace std;

int countRotations(vector<int>& arr) {
    int n = arr.size();
    int lo = 0, hi = n - 1;

    while (lo < hi) {
        if (arr[lo] <= arr[hi]) break; // already sorted

        int mid = lo + (hi - lo) / 2;
        if (arr[mid] >= arr[lo]) {
            lo = mid + 1;
        } else {
            hi = mid;
        }
    }
    return lo; // index of the minimum = number of rotations
}

int main() {
    vector<int> arr = {4, 5, 6, 7, 0, 1, 2};
    cout << "Number of rotations: " << countRotations(arr) << endl;
    return 0;
}
\end{lstlisting}

\begin{complexitybox}
\textbf{Time Complexity.} $O(\log n)$, identical to
Section~\ref{sec:bs-find-min-rotated}.
\textbf{Space Complexity.} $O(1)$ auxiliary space.
\end{complexitybox}

\begin{takeawaybox}
When a new problem statement turns out to want the \emph{index} of a
quantity you already know how to binary search for (rather than its
\emph{value}), there is usually nothing new to derive -- just track the
index alongside, or instead of, the value throughout the same search.
\end{takeawaybox}
